% bismi ALLAH ARRAHMAN ARRAHIM











































\documentclass{article}

\begin{document}

\title{Lois expérimentales révélées par les expériences PCH :
croisement des régimes dimension-calibration et implications pour les LLM}

\maketitle


\section{Résumé expérimental global}

L'ensemble des expériences montre que le PCH (Projection Contextuelle Harmonisée)
ne se comporte pas comme un simple module d'amélioration de features.

Il agit comme un opérateur géométrique contraint :

\[
z \rightarrow P_c(z)
\]

où $P_c$ dépend du contexte latent $c$.

Les expériences historiques
(cas3\_norm\_all\_exps\_0$-$4)
et les nouvelles expériences
(ex14, ex15, ex17, ex18, ex19)
révèlent des lois invariantes.


\section{Loi 1 : invariance par normalisation}

Toutes les expériences montrent :

\[
\gamma \approx 1
\]

\[
\tau \approx 3.4
\]

et dans le régime dimensionnel stable :

\[
\rho_{obs}\approx0.2272727
\]


Cette convergence est indépendante de :

\begin{itemize}
\item $\lambda_{ridge}$
\item ridge ou absence de ridge
\item nombre d'itérations
\item dimension latente modérée
\end{itemize}


Donc :

\[
\boxed{
La normalisation impose un attracteur géométrique universel
}
\]


Le système ne cherche plus le meilleur hyperparamètre.

Il cherche une géométrie satisfaisant les contraintes.


\section{Loi 2 : stabilité dynamique}

Toutes les expériences donnent :

\[
N_{violations}=0
\]

pour la monotonie énergétique.


Le résidu point fixe :

\[
||x-T(x)||
\]

reste :

\[
10^{-17}
\]

à

\[
10^{-16}
\]


Donc :

\[
E(T(x))\leq E(x)
\]


et :

\[
x^*=T(x^*)
\]


Le PCH est un opérateur contractant.


\section{Loi 3 : compression géométrique}

Avant projection :

\[
d_{avant}
\]

Après :

\[
d_{apres}
\]


On observe :

\[
\frac{d_{avant}-d_{apres}}
{d_{avant}}
\approx90\%-95\%
\]


Exemples :

\[
0.036\rightarrow0.0028
\]

(ex14)

\[
0.232\rightarrow0.027
\]

(ex18)


Donc :

\[
\boxed{
Le PCH réduit fortement la distance au contexte latent
}
\]


Ce n'est pas une amélioration statistique classique.

C'est une contraction vers une variété.


\section{Loi 4 : découverte du régime dimension-calibration}

La nouvelle observation fondamentale vient du croisement :

\[
d = \text{latent\_dim}
\]

et :

\[
n_c=n_{calib\_per\_context}
\]


Trois régimes apparaissent.


\subsection{Régime équilibré}

Lorsque :

\[
d=n_c
\]


Exemples :

\[
d=4,\ n_c=4
\]

(ex14)


Le système donne :

\[
\rho=0.22727
\]

\[
operator\ gap\approx10^{-12}
\]

\[
projector\ gap\approx10^{-12}
\]


La projection est quasi exacte.


Conclusion :

\[
\boxed{
d=n_c
\Rightarrow
projection idéale
}
\]


Le nombre de degrés de liberté correspond exactement
au nombre de contraintes contextuelles.


\subsection{Régime sous-calibré}

Lorsque :

\[
d<n_c
\]


Exemple :

\[
d=1,\ n_c=1
\]


(ex18)


On obtient :

\[
\rho=0.089
\]


Le système reste stable :

\[
fixed\ point\approx10^{-17}
\]


mais la capacité contextuelle chute.


La projection devient :

\[
P_c:\mathbb{R}^D\rightarrow\mathbb{R}
\]


Elle détruit trop d'information.


Conséquence :

oracle :

\[
accuracy=1
\]


mais routed :

\[
accuracy<baseline
\]


Donc :

\[
\boxed{
Une projection trop basse dimension devient dangereuse sans contexte parfait
}
\]


\subsection{Régime surdimensionné}

Lorsque :

\[
d>n_c
\]


Exemple :

\[
d=32,\ n_c=48
\]


Ici :

\[
d<n_c
\]

mais le nombre de points de calibration est grand.


On observe :

\[
operator\ gap=10^{-4}
\]

à

\[
10^{-2}
\]


Donc la projection existe mais perd son exactitude.


Le système reste stable mais la géométrie est moins propre.


\section{Nouvelle loi fondamentale}

Le facteur critique n'est pas seulement :

\[
d
\]


mais :

\[
\frac{d}{n_c}
\]


On observe :

\[
\boxed{
\text{qualité géométrique}
=
f(\frac{d}{n_c})
}
\]


avec une zone optimale proche de :

\[
\frac{d}{n_c}\approx1
\]


Trop faible :

\[
\frac{d}{n_c}<1
\]

perte d'information.


Trop grand :

\[
\frac{d}{n_c}>1
\]

sous-détermination géométrique.


\section{Loi 5 : oracle vs routage}

Toutes les expériences confirment :

Oracle :

\[
accuracy\rightarrow1
\]


mais routé :

\[
accuracy\approx0.4-0.8
\]


Donc :

\[
\boxed{
La projection est rarement le problème.
L'identification du contexte est le goulot.
}
\]


Le système sait corriger.

Il ne sait pas toujours choisir la bonne correction.


\section{Conséquence pour les LLM}

Un LLM classique produit :

\[
z=f(x)
\]


Le PCH introduit :

\[
z'=P_c(f(x))
\]


où :

\[
c=\text{contexte latent}
\]


Cela permet :


\subsection{RAG contextuel}

Au lieu de seulement récupérer :

\[
documents\rightarrow embedding
\]


on applique :

\[
embedding\rightarrow P_c(embedding)
\]


Le contexte devient une transformation géométrique.


\subsection{Adaptation multi-domaines}

Un même modèle peut avoir :

\[
P_{medical}
\]

\[
P_{legal}
\]

\[
P_{technical}
\]


sans réentraîner le modèle.


\subsection{Réduction du drift}

Le drift devient :

\[
z_t\rightarrow P_{c_t}(z_t)
\]


La représentation est ramenée dans la variété attendue.


\subsection{Mémoire contextuelle}

La mémoire n'est plus seulement :

\[
store(text)
\]


mais :

\[
store(P_c(z))
\]


Les souvenirs sont normalisés par contexte.


\section{Usage opérationnel innovant}

\begin{enumerate}

\item Adaptateur latent sans fine-tuning :

\[
LLM+P_c
\]

au lieu de LoRA.


\item Correction dynamique d'embeddings.


\item Routage de modèles :

le contexte choisit la projection.


\item Compression intelligente :

garder uniquement les directions utiles.


\item Détection d'anomalies :

distance :

\[
||z-P_c(z)||
\]

devient un score de cohérence.


\item Séparation de domaines dans un même espace latent.

\end{enumerate}


\section{Conclusion finale}

Les expériences révèlent trois lois principales :

\[
\boxed{
La normalisation crée un attracteur universel
}
\]


\[
\boxed{
Le PCH est un opérateur contractant de renormalisation
}
\]


\[
\boxed{
La performance dépend du rapport dimension/calibration
}
\]


La nouvelle variable fondamentale n'est pas :

\[
dim
\]

seul.


Mais :

\[
\boxed{
\frac{dim}{n_{calib\_per\_context}}
}
\]


Le futur axe d'amélioration n'est donc pas seulement une meilleure projection.

C'est un meilleur système d'inférence du contexte.

\end{document}